%% ECON 282E -- Session 5, Deck B3: Chebyshev, and Smolyak as a Basis
%% Source: Quant_Macro Lec.9 frames 54-64 (L1389-1608).  Follows
%% Fernandez-Villaverde & Guerron-Quintana, Solution Methods IV pp.13-21;
%% rewritten here and cited.  One author order corrected.
%% Numbers from tools/figures/s05_numbers.json, key "chebyshev".
%% SMOLYAK IS NOT RECOMPUTED: the point counts come from S04_B1's measured
%% table, key "smolyak" in tools/figures/s04_numbers.json, per slides/SOURCES.md.

%% ECON 282E -- Foundations of Macroeconomics (UC Riverside, Fall 2026)
%% Shared Beamer preamble for every deck in slides/S01 ... slides/S10.
%%
%% Usage (a deck lives in slides/S0X/ and is compiled from that folder):
%%   %% ECON 282E -- Foundations of Macroeconomics (UC Riverside, Fall 2026)
%% Shared Beamer preamble for every deck in slides/S01 ... slides/S10.
%%
%% Usage (a deck lives in slides/S0X/ and is compiled from that folder):
%%   %% ECON 282E -- Foundations of Macroeconomics (UC Riverside, Fall 2026)
%% Shared Beamer preamble for every deck in slides/S01 ... slides/S10.
%%
%% Usage (a deck lives in slides/S0X/ and is compiled from that folder):
%%   \input{../preamble/course_preamble.tex}
%%   \decktitle{1}{A1}{Who I Am \& Why This Course}{September 24, 2026}
%%   \begin{document}
%%   \makedecktitle
%%   ...
%%
%% Adapted from Courses/AI-ECON-2026/source/shared_preamble.tex (Metropolis theme,
%% concept/caution/example/takeaway boxes, \sectiondivider). Changes for this course:
%%   * course title block (\decktitle / \makedecktitle) instead of \makebeamertitle
%%   * \zh{...} typesets Chinese (book titles) under pdflatex via CJKutf8
%%   * researchbox for the "Research in the wild" frames
%%   * section pages off; TikZ libraries and the course map loaded here
%% Build with pdflatex only (two passes); tools/build_slides.py does this for you.

\documentclass[10pt,english,aspectratio=169]{beamer}

\usepackage[T1]{fontenc}
\usepackage{textcomp}
\usepackage[utf8]{inputenc}
\usepackage{babel}
\usepackage{CJKutf8}

\usepackage{amsmath, amssymb, amsthm, graphicx}
\usepackage{booktabs, multirow, tabularx}
\usepackage{tikz}
\usetikzlibrary{positioning,arrows.meta,shapes,calc,fit,backgrounds}
\usepackage{pgfplots}
\pgfplotsset{compat=1.16}
\usepackage[most]{tcolorbox}
\tcbuselibrary{skins, breakable}
\usepackage{listings}
\usepackage{xcolor}

\lstset{
  basicstyle=\ttfamily\small,
  keywordstyle=\color{blue!80!black}\bfseries,
  commentstyle=\color{green!50!black}\itshape,
  stringstyle=\color{orange!80!black},
  backgroundcolor=\color{gray!8},
  frame=single,
  rulecolor=\color{gray!40},
  breaklines=true,
  showstringspaces=false,
  numbers=none,
}

\usetheme{metropolis}
\usefonttheme{professionalfonts}
\metroset{sectionpage=none}

\definecolor{LightGray}{RGB}{230,230,230}
\definecolor{RedTitle}{RGB}{0,0,200}
\definecolor{VeryLightGray}{RGB}{245,245,245}
\definecolor{DarkText}{RGB}{0,0,0}
\definecolor{UCRBlue}{RGB}{0,45,106}
\definecolor{UCRGold}{RGB}{241,171,0}

% Stage colours of the course arc (used by the course map and elsewhere)
\colorlet{StageTrunk}{blue!12}
\colorlet{StageBranch}{cyan!14}
\colorlet{StageMachine}{gray!18}
\colorlet{StageWall}{red!14}
\colorlet{StageTools}{orange!20}
\colorlet{StageData}{green!16}
\colorlet{StageAgents}{violet!14}

\hypersetup{bookmarks=false,breaklinks=true,pdfborder={0 0 0},colorlinks=true,
  linkcolor=DarkText,urlcolor=RedTitle}

\setbeamercolor{background canvas}{bg=white}
\setbeamercolor{title}{fg=RedTitle}
\setbeamercolor{frametitle}{bg=VeryLightGray,fg=DarkText}
\setbeamercolor{normal text}{fg=DarkText}
\setbeamercolor{alerted text}{fg=RedTitle}
\setbeamersize{text margin left=5mm, text margin right=5mm}
\addtobeamertemplate{frametitle}{}{\vspace{-0.5em}}

\setbeamertemplate{navigation symbols}{}
\setbeamertemplate{footline}{%
  \hspace{5mm}{\usebeamerfont{page number in head/foot}\color{black!45}%
  ECON 282E \,$\cdot$\, \insertshortdeck}\hfill%
  \usebeamercolor[fg]{page number in head/foot}%
  \usebeamerfont{page number in head/foot}%
  \insertframenumber\,/\,\inserttotalframenumber\kern1em\vskip2pt%
}

% ---------------------------------------------------------------------------
% Chinese text under pdflatex (book titles etc.)
\newcommand{\zh}[1]{\begin{CJK*}{UTF8}{gbsn}#1\end{CJK*}}

% ---------------------------------------------------------------------------
% Deck title block
%   \decktitle{session}{deck id}{title}{date}
\newcommand{\insertshortdeck}{}
\newcommand{\decksession}{}
\newcommand{\deckid}{}
\newcommand{\decktitletext}{}
\newcommand{\deckdate}{}
\newcommand{\decktitle}[4]{%
  \renewcommand{\decksession}{#1}%
  \renewcommand{\deckid}{#2}%
  \renewcommand{\decktitletext}{#3}%
  \renewcommand{\deckdate}{#4}%
  \renewcommand{\insertshortdeck}{Session #1 \,$\cdot$\, Deck #1.#2}%
  \title{#3}%
  \author{Zhigang Feng}%
  \date{#4}%
}
\newcommand{\makedecktitle}{%
  \begin{frame}[plain,noframenumbering]
    \begin{tikzpicture}[remember picture,overlay]
      \fill[UCRBlue] (current page.north west) rectangle ([yshift=-0.55cm]current page.north east);
      \fill[UCRGold] ([yshift=-0.55cm]current page.north west) rectangle ([yshift=-0.62cm]current page.north east);
      \node[anchor=west, text=white, font=\small\bfseries] at ([xshift=5mm,yshift=-0.275cm]current page.north west)
        {ECON 282E \,$\cdot$\, Foundations of Macroeconomics};
      \node[anchor=east, text=white, font=\small] at ([xshift=-5mm,yshift=-0.275cm]current page.north east)
        {UC Riverside \,$\cdot$\, Fall 2026};
    \end{tikzpicture}
    \vspace*{1.1cm}
    {\usebeamerfont{title}\color{black!55}\large Session \decksession\ \,$\cdot$\, Deck \decksession.\deckid\par}
    \vspace{0.25cm}
    {\usebeamerfont{title}\color{RedTitle}\LARGE\bfseries \decktitletext\par}
    \vspace{0.7cm}
    {\large Zhigang Feng\par}
    \vspace{0.15cm}
    {\color{black!65}\deckdate\par}
    \vfill
    {\footnotesize\itshape\color{black!55}Build it on paper $\;\to\;$ solve it on the machine $\;\to\;$ push it to the frontier with AI\par}
  \end{frame}}

% ---------------------------------------------------------------------------
% Section divider (plain frame, not counted)
\newcommand{\sectiondivider}[2]{%
\begin{frame}[plain,noframenumbering]
\begin{center}
\vfill
{\Large\textbf{#1}\par}
\vspace{0.35cm}
{\normalsize #2\par}
\vfill
\end{center}
\end{frame}
}

% ---------------------------------------------------------------------------
% Boxes
\newtcolorbox{conceptbox}[1][]{colback=blue!3, colframe=RedTitle!55, boxrule=0.35mm,
  arc=1mm, left=2mm, right=2mm, top=1mm, bottom=1mm, #1}
\newtcolorbox{cautionbox}[1][]{colback=red!3, colframe=red!50!black, boxrule=0.35mm,
  arc=1mm, left=2mm, right=2mm, top=1mm, bottom=1mm, #1}
\newtcolorbox{examplebox}[1][]{colback=green!3, colframe=green!45!black, boxrule=0.35mm,
  arc=1mm, left=2mm, right=2mm, top=1mm, bottom=1mm, #1}
\newtcolorbox{takeawaybox}[1][]{colback=VeryLightGray, colframe=RedTitle!55, boxrule=0.35mm,
  arc=1mm, left=2mm, right=2mm, top=1mm, bottom=1mm, #1}
\newtcolorbox{researchbox}[1][]{enhanced, colback=violet!4, colframe=violet!60!black,
  boxrule=0.35mm, arc=1mm, left=2mm, right=2mm, top=1mm, bottom=1mm,
  fonttitle=\bfseries\small, title={Research in the wild}, #1}

\newcommand{\compacttables}{\renewcommand{\arraystretch}{1.15}}
\newcommand{\src}[1]{{\tiny\color{black!55}#1}}

% ---------------------------------------------------------------------------
\input{../preamble/course_map.tex}

%%   \decktitle{1}{A1}{Who I Am \& Why This Course}{September 24, 2026}
%%   \begin{document}
%%   \makedecktitle
%%   ...
%%
%% Adapted from Courses/AI-ECON-2026/source/shared_preamble.tex (Metropolis theme,
%% concept/caution/example/takeaway boxes, \sectiondivider). Changes for this course:
%%   * course title block (\decktitle / \makedecktitle) instead of \makebeamertitle
%%   * \zh{...} typesets Chinese (book titles) under pdflatex via CJKutf8
%%   * researchbox for the "Research in the wild" frames
%%   * section pages off; TikZ libraries and the course map loaded here
%% Build with pdflatex only (two passes); tools/build_slides.py does this for you.

\documentclass[10pt,english,aspectratio=169]{beamer}

\usepackage[T1]{fontenc}
\usepackage{textcomp}
\usepackage[utf8]{inputenc}
\usepackage{babel}
\usepackage{CJKutf8}

\usepackage{amsmath, amssymb, amsthm, graphicx}
\usepackage{booktabs, multirow, tabularx}
\usepackage{tikz}
\usetikzlibrary{positioning,arrows.meta,shapes,calc,fit,backgrounds}
\usepackage{pgfplots}
\pgfplotsset{compat=1.16}
\usepackage[most]{tcolorbox}
\tcbuselibrary{skins, breakable}
\usepackage{listings}
\usepackage{xcolor}

\lstset{
  basicstyle=\ttfamily\small,
  keywordstyle=\color{blue!80!black}\bfseries,
  commentstyle=\color{green!50!black}\itshape,
  stringstyle=\color{orange!80!black},
  backgroundcolor=\color{gray!8},
  frame=single,
  rulecolor=\color{gray!40},
  breaklines=true,
  showstringspaces=false,
  numbers=none,
}

\usetheme{metropolis}
\usefonttheme{professionalfonts}
\metroset{sectionpage=none}

\definecolor{LightGray}{RGB}{230,230,230}
\definecolor{RedTitle}{RGB}{0,0,200}
\definecolor{VeryLightGray}{RGB}{245,245,245}
\definecolor{DarkText}{RGB}{0,0,0}
\definecolor{UCRBlue}{RGB}{0,45,106}
\definecolor{UCRGold}{RGB}{241,171,0}

% Stage colours of the course arc (used by the course map and elsewhere)
\colorlet{StageTrunk}{blue!12}
\colorlet{StageBranch}{cyan!14}
\colorlet{StageMachine}{gray!18}
\colorlet{StageWall}{red!14}
\colorlet{StageTools}{orange!20}
\colorlet{StageData}{green!16}
\colorlet{StageAgents}{violet!14}

\hypersetup{bookmarks=false,breaklinks=true,pdfborder={0 0 0},colorlinks=true,
  linkcolor=DarkText,urlcolor=RedTitle}

\setbeamercolor{background canvas}{bg=white}
\setbeamercolor{title}{fg=RedTitle}
\setbeamercolor{frametitle}{bg=VeryLightGray,fg=DarkText}
\setbeamercolor{normal text}{fg=DarkText}
\setbeamercolor{alerted text}{fg=RedTitle}
\setbeamersize{text margin left=5mm, text margin right=5mm}
\addtobeamertemplate{frametitle}{}{\vspace{-0.5em}}

\setbeamertemplate{navigation symbols}{}
\setbeamertemplate{footline}{%
  \hspace{5mm}{\usebeamerfont{page number in head/foot}\color{black!45}%
  ECON 282E \,$\cdot$\, \insertshortdeck}\hfill%
  \usebeamercolor[fg]{page number in head/foot}%
  \usebeamerfont{page number in head/foot}%
  \insertframenumber\,/\,\inserttotalframenumber\kern1em\vskip2pt%
}

% ---------------------------------------------------------------------------
% Chinese text under pdflatex (book titles etc.)
\newcommand{\zh}[1]{\begin{CJK*}{UTF8}{gbsn}#1\end{CJK*}}

% ---------------------------------------------------------------------------
% Deck title block
%   \decktitle{session}{deck id}{title}{date}
\newcommand{\insertshortdeck}{}
\newcommand{\decksession}{}
\newcommand{\deckid}{}
\newcommand{\decktitletext}{}
\newcommand{\deckdate}{}
\newcommand{\decktitle}[4]{%
  \renewcommand{\decksession}{#1}%
  \renewcommand{\deckid}{#2}%
  \renewcommand{\decktitletext}{#3}%
  \renewcommand{\deckdate}{#4}%
  \renewcommand{\insertshortdeck}{Session #1 \,$\cdot$\, Deck #1.#2}%
  \title{#3}%
  \author{Zhigang Feng}%
  \date{#4}%
}
\newcommand{\makedecktitle}{%
  \begin{frame}[plain,noframenumbering]
    \begin{tikzpicture}[remember picture,overlay]
      \fill[UCRBlue] (current page.north west) rectangle ([yshift=-0.55cm]current page.north east);
      \fill[UCRGold] ([yshift=-0.55cm]current page.north west) rectangle ([yshift=-0.62cm]current page.north east);
      \node[anchor=west, text=white, font=\small\bfseries] at ([xshift=5mm,yshift=-0.275cm]current page.north west)
        {ECON 282E \,$\cdot$\, Foundations of Macroeconomics};
      \node[anchor=east, text=white, font=\small] at ([xshift=-5mm,yshift=-0.275cm]current page.north east)
        {UC Riverside \,$\cdot$\, Fall 2026};
    \end{tikzpicture}
    \vspace*{1.1cm}
    {\usebeamerfont{title}\color{black!55}\large Session \decksession\ \,$\cdot$\, Deck \decksession.\deckid\par}
    \vspace{0.25cm}
    {\usebeamerfont{title}\color{RedTitle}\LARGE\bfseries \decktitletext\par}
    \vspace{0.7cm}
    {\large Zhigang Feng\par}
    \vspace{0.15cm}
    {\color{black!65}\deckdate\par}
    \vfill
    {\footnotesize\itshape\color{black!55}Build it on paper $\;\to\;$ solve it on the machine $\;\to\;$ push it to the frontier with AI\par}
  \end{frame}}

% ---------------------------------------------------------------------------
% Section divider (plain frame, not counted)
\newcommand{\sectiondivider}[2]{%
\begin{frame}[plain,noframenumbering]
\begin{center}
\vfill
{\Large\textbf{#1}\par}
\vspace{0.35cm}
{\normalsize #2\par}
\vfill
\end{center}
\end{frame}
}

% ---------------------------------------------------------------------------
% Boxes
\newtcolorbox{conceptbox}[1][]{colback=blue!3, colframe=RedTitle!55, boxrule=0.35mm,
  arc=1mm, left=2mm, right=2mm, top=1mm, bottom=1mm, #1}
\newtcolorbox{cautionbox}[1][]{colback=red!3, colframe=red!50!black, boxrule=0.35mm,
  arc=1mm, left=2mm, right=2mm, top=1mm, bottom=1mm, #1}
\newtcolorbox{examplebox}[1][]{colback=green!3, colframe=green!45!black, boxrule=0.35mm,
  arc=1mm, left=2mm, right=2mm, top=1mm, bottom=1mm, #1}
\newtcolorbox{takeawaybox}[1][]{colback=VeryLightGray, colframe=RedTitle!55, boxrule=0.35mm,
  arc=1mm, left=2mm, right=2mm, top=1mm, bottom=1mm, #1}
\newtcolorbox{researchbox}[1][]{enhanced, colback=violet!4, colframe=violet!60!black,
  boxrule=0.35mm, arc=1mm, left=2mm, right=2mm, top=1mm, bottom=1mm,
  fonttitle=\bfseries\small, title={Research in the wild}, #1}

\newcommand{\compacttables}{\renewcommand{\arraystretch}{1.15}}
\newcommand{\src}[1]{{\tiny\color{black!55}#1}}

% ---------------------------------------------------------------------------
%% Course map: the recurring "you are here" slide for ECON 282E.
%% Loaded by course_preamble.tex. Pattern follows AI-ECON-2026 lec01_map.tex, but the map
%% shows the whole ten-session course rather than one lecture.
%%
%%   \CourseMap{s}{label}   s = 1..10 highlights session s and prints "you are here: label"
%%                          (e.g. \CourseMap{1}{1.A2}); earlier sessions get a check mark.
%%   \CourseMap{0}{}        full overview, nothing highlighted.
%%
%% Note: TikZ option lists are comma-split before TeX conditionals run, so every \ifnum
%% below sits outside [...] option lists.

\newcommand{\CourseMap}[2]{%
\begin{center}
\begin{tikzpicture}[
    sess/.style={rectangle, rounded corners=2pt, draw=black!45, minimum width=1.34cm,
        minimum height=1.12cm, text width=1.26cm, align=center, inner sep=1pt},
    stage/.style={font=\tiny\bfseries, text=black!70},
]
  \foreach \i/\d/\t/\c in {%
      1/Sep 24/Intro \\ Growth/StageTrunk,
      2/Oct 1/HA $\cdot$ OLG \\ Policy/StageBranch,
      3/Oct 8/Num.\ DP \\ Python/StageMachine,
      4/Oct 15/Numerics \\ I/StageMachine,
      5/Oct 22/Numerics \\ II/StageMachine,
      6/Oct 29/The wall \\ \textbf{Midterm}/StageWall,
      7/Nov 5/ML \\ Deep nets/StageTools,
      8/Nov 12/RL \\ HA $+$ DL/StageTools,
      9/Nov 19/Text \\ LLMs/StageData,
      10/Dec 3/Agents \\ \textbf{Projects}/StageAgents} {
    \node[sess, fill=\c] (n\i) at ({1.46*(\i-1)}, 0)
      {{\scriptsize\bfseries S\i}\\[-1pt]{\tiny\color{black!60}\d}\\[1pt]{\tiny \t}};
    \ifnum\i<#1
      \node[font=\scriptsize, text=green!45!black] at ([xshift=-2.5mm,yshift=-2.2mm]n\i.north east) {$\checkmark$};
    \fi
    \ifnum\i=#1
      \draw[RedTitle, very thick, rounded corners=2pt]
        ([xshift=-0.6mm,yshift=-0.6mm]n\i.south west) rectangle ([xshift=0.6mm,yshift=0.6mm]n\i.north east);
      % keep the label inside the picture at both ends of the row
      \ifnum\i<3
        \node[font=\scriptsize\bfseries, text=RedTitle, anchor=south west] at ([yshift=1.2mm]n\i.north west)
          {$\blacktriangledown$\,you are here: #2};
      \else\ifnum\i>8
        \node[font=\scriptsize\bfseries, text=RedTitle, anchor=south east] at ([yshift=1.2mm]n\i.north east)
          {you are here: #2\,$\blacktriangledown$};
      \else
        \node[font=\scriptsize\bfseries, text=RedTitle, anchor=south] at ([yshift=1.2mm]n\i.north)
          {you are here: #2\,$\blacktriangledown$};
      \fi\fi
    \fi
  }
  % stage band under the sessions
  \foreach \a/\b/\lab/\c in {1/1/trunk/StageTrunk, 2/2/branches/StageBranch,
      3/5/the machine $\cdot$ classical methods/StageMachine, 6/6/the wall/StageWall,
      7/8/new tools/StageTools, 9/9/new data/StageData, 10/10/colleagues/StageAgents} {
    \fill[\c] ([yshift=-1.5mm]n\a.south west) rectangle ([yshift=-4.6mm]n\b.south east);
    \path ([yshift=-1.5mm]n\a.south west) -- ([yshift=-4.6mm]n\b.south east)
      node[midway, stage] {\lab};
  }
  \node[font=\footnotesize\itshape, text=black!70] at ([yshift=-9.5mm]$(n1.south)!0.5!(n10.south)$)
    {Build it on paper $\;\to\;$ solve it on the machine $\;\to\;$ push it to the frontier with AI};
\end{tikzpicture}
\end{center}
}


%%   \decktitle{1}{A1}{Who I Am \& Why This Course}{September 24, 2026}
%%   \begin{document}
%%   \makedecktitle
%%   ...
%%
%% Adapted from Courses/AI-ECON-2026/source/shared_preamble.tex (Metropolis theme,
%% concept/caution/example/takeaway boxes, \sectiondivider). Changes for this course:
%%   * course title block (\decktitle / \makedecktitle) instead of \makebeamertitle
%%   * \zh{...} typesets Chinese (book titles) under pdflatex via CJKutf8
%%   * researchbox for the "Research in the wild" frames
%%   * section pages off; TikZ libraries and the course map loaded here
%% Build with pdflatex only (two passes); tools/build_slides.py does this for you.

\documentclass[10pt,english,aspectratio=169]{beamer}

\usepackage[T1]{fontenc}
\usepackage{textcomp}
\usepackage[utf8]{inputenc}
\usepackage{babel}
\usepackage{CJKutf8}

\usepackage{amsmath, amssymb, amsthm, graphicx}
\usepackage{booktabs, multirow, tabularx}
\usepackage{tikz}
\usetikzlibrary{positioning,arrows.meta,shapes,calc,fit,backgrounds}
\usepackage{pgfplots}
\pgfplotsset{compat=1.16}
\usepackage[most]{tcolorbox}
\tcbuselibrary{skins, breakable}
\usepackage{listings}
\usepackage{xcolor}

\lstset{
  basicstyle=\ttfamily\small,
  keywordstyle=\color{blue!80!black}\bfseries,
  commentstyle=\color{green!50!black}\itshape,
  stringstyle=\color{orange!80!black},
  backgroundcolor=\color{gray!8},
  frame=single,
  rulecolor=\color{gray!40},
  breaklines=true,
  showstringspaces=false,
  numbers=none,
}

\usetheme{metropolis}
\usefonttheme{professionalfonts}
\metroset{sectionpage=none}

\definecolor{LightGray}{RGB}{230,230,230}
\definecolor{RedTitle}{RGB}{0,0,200}
\definecolor{VeryLightGray}{RGB}{245,245,245}
\definecolor{DarkText}{RGB}{0,0,0}
\definecolor{UCRBlue}{RGB}{0,45,106}
\definecolor{UCRGold}{RGB}{241,171,0}

% Stage colours of the course arc (used by the course map and elsewhere)
\colorlet{StageTrunk}{blue!12}
\colorlet{StageBranch}{cyan!14}
\colorlet{StageMachine}{gray!18}
\colorlet{StageWall}{red!14}
\colorlet{StageTools}{orange!20}
\colorlet{StageData}{green!16}
\colorlet{StageAgents}{violet!14}

\hypersetup{bookmarks=false,breaklinks=true,pdfborder={0 0 0},colorlinks=true,
  linkcolor=DarkText,urlcolor=RedTitle}

\setbeamercolor{background canvas}{bg=white}
\setbeamercolor{title}{fg=RedTitle}
\setbeamercolor{frametitle}{bg=VeryLightGray,fg=DarkText}
\setbeamercolor{normal text}{fg=DarkText}
\setbeamercolor{alerted text}{fg=RedTitle}
\setbeamersize{text margin left=5mm, text margin right=5mm}
\addtobeamertemplate{frametitle}{}{\vspace{-0.5em}}

\setbeamertemplate{navigation symbols}{}
\setbeamertemplate{footline}{%
  \hspace{5mm}{\usebeamerfont{page number in head/foot}\color{black!45}%
  ECON 282E \,$\cdot$\, \insertshortdeck}\hfill%
  \usebeamercolor[fg]{page number in head/foot}%
  \usebeamerfont{page number in head/foot}%
  \insertframenumber\,/\,\inserttotalframenumber\kern1em\vskip2pt%
}

% ---------------------------------------------------------------------------
% Chinese text under pdflatex (book titles etc.)
\newcommand{\zh}[1]{\begin{CJK*}{UTF8}{gbsn}#1\end{CJK*}}

% ---------------------------------------------------------------------------
% Deck title block
%   \decktitle{session}{deck id}{title}{date}
\newcommand{\insertshortdeck}{}
\newcommand{\decksession}{}
\newcommand{\deckid}{}
\newcommand{\decktitletext}{}
\newcommand{\deckdate}{}
\newcommand{\decktitle}[4]{%
  \renewcommand{\decksession}{#1}%
  \renewcommand{\deckid}{#2}%
  \renewcommand{\decktitletext}{#3}%
  \renewcommand{\deckdate}{#4}%
  \renewcommand{\insertshortdeck}{Session #1 \,$\cdot$\, Deck #1.#2}%
  \title{#3}%
  \author{Zhigang Feng}%
  \date{#4}%
}
\newcommand{\makedecktitle}{%
  \begin{frame}[plain,noframenumbering]
    \begin{tikzpicture}[remember picture,overlay]
      \fill[UCRBlue] (current page.north west) rectangle ([yshift=-0.55cm]current page.north east);
      \fill[UCRGold] ([yshift=-0.55cm]current page.north west) rectangle ([yshift=-0.62cm]current page.north east);
      \node[anchor=west, text=white, font=\small\bfseries] at ([xshift=5mm,yshift=-0.275cm]current page.north west)
        {ECON 282E \,$\cdot$\, Foundations of Macroeconomics};
      \node[anchor=east, text=white, font=\small] at ([xshift=-5mm,yshift=-0.275cm]current page.north east)
        {UC Riverside \,$\cdot$\, Fall 2026};
    \end{tikzpicture}
    \vspace*{1.1cm}
    {\usebeamerfont{title}\color{black!55}\large Session \decksession\ \,$\cdot$\, Deck \decksession.\deckid\par}
    \vspace{0.25cm}
    {\usebeamerfont{title}\color{RedTitle}\LARGE\bfseries \decktitletext\par}
    \vspace{0.7cm}
    {\large Zhigang Feng\par}
    \vspace{0.15cm}
    {\color{black!65}\deckdate\par}
    \vfill
    {\footnotesize\itshape\color{black!55}Build it on paper $\;\to\;$ solve it on the machine $\;\to\;$ push it to the frontier with AI\par}
  \end{frame}}

% ---------------------------------------------------------------------------
% Section divider (plain frame, not counted)
\newcommand{\sectiondivider}[2]{%
\begin{frame}[plain,noframenumbering]
\begin{center}
\vfill
{\Large\textbf{#1}\par}
\vspace{0.35cm}
{\normalsize #2\par}
\vfill
\end{center}
\end{frame}
}

% ---------------------------------------------------------------------------
% Boxes
\newtcolorbox{conceptbox}[1][]{colback=blue!3, colframe=RedTitle!55, boxrule=0.35mm,
  arc=1mm, left=2mm, right=2mm, top=1mm, bottom=1mm, #1}
\newtcolorbox{cautionbox}[1][]{colback=red!3, colframe=red!50!black, boxrule=0.35mm,
  arc=1mm, left=2mm, right=2mm, top=1mm, bottom=1mm, #1}
\newtcolorbox{examplebox}[1][]{colback=green!3, colframe=green!45!black, boxrule=0.35mm,
  arc=1mm, left=2mm, right=2mm, top=1mm, bottom=1mm, #1}
\newtcolorbox{takeawaybox}[1][]{colback=VeryLightGray, colframe=RedTitle!55, boxrule=0.35mm,
  arc=1mm, left=2mm, right=2mm, top=1mm, bottom=1mm, #1}
\newtcolorbox{researchbox}[1][]{enhanced, colback=violet!4, colframe=violet!60!black,
  boxrule=0.35mm, arc=1mm, left=2mm, right=2mm, top=1mm, bottom=1mm,
  fonttitle=\bfseries\small, title={Research in the wild}, #1}

\newcommand{\compacttables}{\renewcommand{\arraystretch}{1.15}}
\newcommand{\src}[1]{{\tiny\color{black!55}#1}}

% ---------------------------------------------------------------------------
%% Course map: the recurring "you are here" slide for ECON 282E.
%% Loaded by course_preamble.tex. Pattern follows AI-ECON-2026 lec01_map.tex, but the map
%% shows the whole ten-session course rather than one lecture.
%%
%%   \CourseMap{s}{label}   s = 1..10 highlights session s and prints "you are here: label"
%%                          (e.g. \CourseMap{1}{1.A2}); earlier sessions get a check mark.
%%   \CourseMap{0}{}        full overview, nothing highlighted.
%%
%% Note: TikZ option lists are comma-split before TeX conditionals run, so every \ifnum
%% below sits outside [...] option lists.

\newcommand{\CourseMap}[2]{%
\begin{center}
\begin{tikzpicture}[
    sess/.style={rectangle, rounded corners=2pt, draw=black!45, minimum width=1.34cm,
        minimum height=1.12cm, text width=1.26cm, align=center, inner sep=1pt},
    stage/.style={font=\tiny\bfseries, text=black!70},
]
  \foreach \i/\d/\t/\c in {%
      1/Sep 24/Intro \\ Growth/StageTrunk,
      2/Oct 1/HA $\cdot$ OLG \\ Policy/StageBranch,
      3/Oct 8/Num.\ DP \\ Python/StageMachine,
      4/Oct 15/Numerics \\ I/StageMachine,
      5/Oct 22/Numerics \\ II/StageMachine,
      6/Oct 29/The wall \\ \textbf{Midterm}/StageWall,
      7/Nov 5/ML \\ Deep nets/StageTools,
      8/Nov 12/RL \\ HA $+$ DL/StageTools,
      9/Nov 19/Text \\ LLMs/StageData,
      10/Dec 3/Agents \\ \textbf{Projects}/StageAgents} {
    \node[sess, fill=\c] (n\i) at ({1.46*(\i-1)}, 0)
      {{\scriptsize\bfseries S\i}\\[-1pt]{\tiny\color{black!60}\d}\\[1pt]{\tiny \t}};
    \ifnum\i<#1
      \node[font=\scriptsize, text=green!45!black] at ([xshift=-2.5mm,yshift=-2.2mm]n\i.north east) {$\checkmark$};
    \fi
    \ifnum\i=#1
      \draw[RedTitle, very thick, rounded corners=2pt]
        ([xshift=-0.6mm,yshift=-0.6mm]n\i.south west) rectangle ([xshift=0.6mm,yshift=0.6mm]n\i.north east);
      % keep the label inside the picture at both ends of the row
      \ifnum\i<3
        \node[font=\scriptsize\bfseries, text=RedTitle, anchor=south west] at ([yshift=1.2mm]n\i.north west)
          {$\blacktriangledown$\,you are here: #2};
      \else\ifnum\i>8
        \node[font=\scriptsize\bfseries, text=RedTitle, anchor=south east] at ([yshift=1.2mm]n\i.north east)
          {you are here: #2\,$\blacktriangledown$};
      \else
        \node[font=\scriptsize\bfseries, text=RedTitle, anchor=south] at ([yshift=1.2mm]n\i.north)
          {you are here: #2\,$\blacktriangledown$};
      \fi\fi
    \fi
  }
  % stage band under the sessions
  \foreach \a/\b/\lab/\c in {1/1/trunk/StageTrunk, 2/2/branches/StageBranch,
      3/5/the machine $\cdot$ classical methods/StageMachine, 6/6/the wall/StageWall,
      7/8/new tools/StageTools, 9/9/new data/StageData, 10/10/colleagues/StageAgents} {
    \fill[\c] ([yshift=-1.5mm]n\a.south west) rectangle ([yshift=-4.6mm]n\b.south east);
    \path ([yshift=-1.5mm]n\a.south west) -- ([yshift=-4.6mm]n\b.south east)
      node[midway, stage] {\lab};
  }
  \node[font=\footnotesize\itshape, text=black!70] at ([yshift=-9.5mm]$(n1.south)!0.5!(n10.south)$)
    {Build it on paper $\;\to\;$ solve it on the machine $\;\to\;$ push it to the frontier with AI};
\end{tikzpicture}
\end{center}
}


\graphicspath{{./figures/}}

\lstset{language=Python, basicstyle=\ttfamily\scriptsize, frame=none,
        backgroundcolor=\color{gray!8}, aboveskip=2pt, belowskip=2pt}

\decktitle{5}{B3}{Chebyshev, and Smolyak as a Basis}{Thursday, October 22, 2026}

\begin{document}

\makedecktitle

%=============================================================================
\begin{frame}{Where We Are}
\CourseMap{5}{5.B3}
\vspace{-1mm}
\begin{center}
\small \textbf{Block B:} 5.B1 the residual $\;\cdot\;$ 5.B2 choosing a basis $\;\cdot\;$ \textbf{5.B3} Chebyshev $\;\cdot\;$ 5.B4 finite elements $\;\cdot\;$ 5.B5 determining the coefficients $\;\cdot\;$ 5.B6 the worked example.
\end{center}
\end{frame}

%=============================================================================
\begin{frame}{One Family, Two Famous Members}
\begin{columns}[T]
\begin{column}{0.55\textwidth}
\small
The \textbf{Jacobi} (or hypergeometric) family $P_n^{a,b}$, for $a,b>-1$, is defined by an orthogonality condition with a weight:
\[
  \int_{-1}^{1}(1-x)^{a}(1+x)^{b}\,P_n^{a,b}(x)P_m^{a,b}(x)\,dx=0,\quad m\neq n .
\]
The weight decides which family you get, and $a,b>-1$ keeps the integral finite.
\medskip
\begin{itemize}\small
  \item $a=b=-\tfrac12$: \textbf{Chebyshev}. The weight blows up at the endpoints, so the family pays extra attention there.
  \item $a=b=0$: \textbf{Legendre}. Uniform weight; the right family for Gauss--Legendre quadrature, which 4.B2 already used.
\end{itemize}
\end{column}
\begin{column}{0.42\textwidth}
\begin{conceptbox}[title=\textbf{Notation}]\small
The literature writes the Jacobi parameters as $\alpha,\beta$. In this course $\alpha$ is the capital share and $\beta$ the discount factor, so here they are $a,b$.
\end{conceptbox}
\medskip
\begin{examplebox}\footnotesize
That endpoint weighting is not a curiosity. It is why Chebyshev nodes cluster at the ends, and why they defeat Runge's phenomenon (4.B1).
\end{examplebox}
\end{column}
\end{columns}
\vspace{1mm}
\src{Quant\_Macro Lec.~9, ``Orthogonal Polynomials''.}
\end{frame}

%=============================================================================
\begin{frame}{Why Chebyshev Is the Default}
\begin{columns}[T]
\begin{column}{0.56\textwidth}
\small
\begin{itemize}\small\setlength{\itemsep}{1.2mm}
  \item Many \textbf{simple closed-form expressions} are available.
  \item More \textbf{robust than the alternatives} for interpolation.
  \item \textbf{Bounded} in $[-1,1]$ --- so no coefficient has to cancel a large number, which is precisely what killed the monomials.
  \item \textbf{Smooth}.
  \item Moving between coefficients and values at the Chebyshev nodes is a \textbf{cosine transform} --- an FFT, $O(n\log n)$.
  \item Several \textbf{theorems bound the interpolation error}, so the choice can be defended rather than asserted.
\end{itemize}
\end{column}
\begin{column}{0.41\textwidth}
\begin{takeawaybox}\footnotesize
Standard references: Boyd (2001), \emph{Chebyshev and Fourier Spectral Methods}; Fornberg (1998), \emph{A Practical Guide to Pseudospectral Methods}.
\end{takeawaybox}
\medskip
\begin{examplebox}\footnotesize
Boyd's own rule of thumb is worth carrying: \textbf{when in doubt, use Chebyshev} --- unless the solution is periodic, in which case use Fourier.
\end{examplebox}
\end{column}
\end{columns}
\vspace{1mm}
\src{Quant\_Macro Lec.~9, ``Chebyshev polynomials'' (advantages frame).}
\end{frame}

%=============================================================================
\begin{frame}{The Recursion and the Nodes}
\begin{columns}[T]
\begin{column}{0.54\textwidth}
\small
\begin{align*}
T_0(x)&=1\\
T_1(x)&=x\\
T_{n+1}(x)&=2x\,T_n(x)-T_{n-1}(x)
\end{align*}
giving $1,\;x,\;2x^{2}-1,\;4x^{3}-3x,\;8x^{4}-8x^{2}+1,\dots$
\medskip

The $n$ roots of $T_n$ are
\[
  x_k=\cos\frac{2k-1}{2n}\pi,\qquad k=1,\dots,n .
\]
They are \textbf{not} equally spaced: they cluster quadratically towards $\pm1$, which is exactly where an equispaced interpolant misbehaves.
\end{column}
\begin{column}{0.43\textwidth}
\begin{conceptbox}[title=\textbf{Where you have seen this}]\small
4.B1 measured it: on Runge's function, equispaced interpolation \emph{diverges} as the degree rises, while Chebyshev nodes converge. Same nodes, and now the polynomials that go with them.
\end{conceptbox}
\medskip
\begin{examplebox}\footnotesize
Three lines of code. The recursion is numerically stable, which is more than can be said for computing $x^{n}$ directly.
\end{examplebox}
\end{column}
\end{columns}
\vspace{1mm}
\src{Quant\_Macro Lec.~9, ``Chebyshev polynomials'' (recursion frame).}
\end{frame}

%=============================================================================
\begin{frame}{Closed Forms, and the Change of Variable}
\begin{columns}[T]
\begin{column}{0.52\textwidth}
\small
\[
T_n(x)=\begin{cases}
\cos\big(n\arccos x\big) & -1\le x\le 1\\
\cosh\big(n\,\mathrm{arcosh}\,x\big) & x\ge 1\\
(-1)^{n}\cosh\big(n\,\mathrm{arcosh}(-x)\big) & x\le -1
\end{cases}
\]
\medskip
A polynomial written as a cosine: that identity is the reason the coefficient transform is an FFT, and the reason $|T_n|\le1$ on the interval.
\medskip

Our state space is not $[-1,1]$, so map it in:
\[
  x \;\longmapsto\; 2\,\frac{x-a}{b-a}-1 .
\]
\end{column}
\begin{column}{0.45\textwidth}
\begin{conceptbox}[title=\textbf{The weight}]\small
Chebyshev polynomials are orthogonal with respect to
\[
  w(x)=\frac{1}{\sqrt{1-x^{2}}} ,
\]
which is infinite at $\pm1$ --- the analytical statement of the same endpoint emphasis.
\end{conceptbox}
\medskip
\begin{cautionbox}\footnotesize
Choosing $[a,b]$ is a modelling decision, not a formality: outside it the polynomial is not just inaccurate, it diverges like $\cosh$.
\end{cautionbox}
\end{column}
\end{columns}
\vspace{1mm}
\src{Quant\_Macro Lec.~9, ``Chebyshev polynomials'' (closed-form and domain frames).}
\end{frame}

%=============================================================================
\begin{frame}{The Interpolation Theorem}
\begin{columns}[T]
\begin{column}{0.55\textwidth}
\small
\begin{conceptbox}[title=\textbf{Chebyshev Interpolation Theorem}]\small
If an approximating polynomial is made exact at the roots of $T_n$, then as $n\to\infty$ the approximation error becomes arbitrarily small --- for any continuous function.
\end{conceptbox}
\medskip
There is a sharper version worth knowing, because it is what makes the choice \emph{defensible} rather than merely traditional. Chebyshev-node interpolation is never much worse than the \textbf{best possible} polynomial of the same degree:
\[
  \|f-p_n\|_{\infty}\;\le\;\Big(\tfrac{2}{\pi}\log(n+1)+1\Big)\;\|f-p_n^{\text{best}}\|_{\infty}.
\]
At $n=33$ that factor is $\mathbf{3.24}$.
\end{column}
\begin{column}{0.42\textwidth}
\begin{takeawaybox}\footnotesize
So you are within a factor of three of \textbf{the best any polynomial of that degree could do} --- and the best one is not computable, while this one is an FFT.
\end{takeawaybox}
\medskip
\begin{cautionbox}\footnotesize
The factor grows like $\log n$, so it is still $4.0$ at $n=1000$. It never bites.
\end{cautionbox}
\end{column}
\end{columns}
\vspace{1mm}
\src{Quant\_Macro Lec.~9, ``Chebyshev polynomials'' (interpolation-theorem frame). Ledger \texttt{chebyshev.lebesgue\_factor\_n33} $=3.2450$.}
\end{frame}

%=============================================================================
\begin{frame}{Measured: What Smoothness Buys, and What a Kink Costs}
\begin{columns}[T]
\begin{column}{0.5\textwidth}
\begin{center}
\begin{tikzpicture}
\begin{axis}[
    width=7.2cm, height=4.4cm,
    xmin=0, xmax=32, ymin=-16, ymax=1,
    xlabel={\scriptsize coefficient index}, ylabel={\scriptsize $\log_{10}|\theta_i|$},
    tick label style={font=\scriptsize}, label style={font=\scriptsize},
    scaled ticks=false, /pgf/number format/fixed,
    /pgf/number format/precision=0,
    axis x line*=bottom, axis y line*=left,
    legend style={font=\tiny, draw=none, fill=none,
                  at={(0.98,0.97)}, anchor=north east}]
  \addplot[very thick, RedTitle] coordinates {
      (0,-0.432) (1,-0.063) (2,-1.295) (3,-0.169) (4,-0.544) (5,-1.581) (6,-1.492) (7,-2.446) 
      (8,-2.982) (9,-3.591) (10,-5.552) (11,-5.200) (12,-6.305) (13,-7.194) (14,-7.933) 
      (15,-10.597) (16,-9.914) (17,-11.121) (18,-12.237) (19,-13.059) (20,-16.535) (21,-16.319) 
      (22,-15.505) (23,-15.422) (24,-15.346) (25,-16.524) (26,-16.003) (27,-15.304) (28,-14.979) 
      (29,-15.700) (30,-16.988) (31,-15.177) (32,-15.217) };
  \addlegendentry{$e^{x}\sin 3x$}
  \addplot[very thick, orange!85!black] coordinates {
      (0,-0.196) (1,-15.696) (2,-0.371) (3,-15.946) (4,-1.066) (5,-16.038) (6,-1.428) (7,-15.659) 
      (8,-1.674) (9,-15.812) (10,-1.858) (11,-15.727) (12,-2.003) (13,-15.920) (14,-2.120) 
      (15,-15.594) (16,-2.216) (17,-16.110) (18,-2.296) (19,-15.567) (20,-2.361) (21,-15.532) 
      (22,-2.414) (23,-15.689) (24,-2.457) (25,-16.291) (26,-2.490) (27,-16.017) (28,-2.515) 
      (29,-15.747) (30,-2.531) (31,-15.897) (32,-2.539) };
  \addlegendentry{$|x|$}
\end{axis}
\end{tikzpicture}
\end{center}
\end{column}
\begin{column}{0.47\textwidth}
\small
\textbf{Each line is one function's list of Chebyshev coefficients} $\theta_i$: $e^{x}\sin 3x$ is
analytic, $|x|$ has a corner mid-interval. Truncating at $n$ drops the tail, so \textbf{the error is
about the first coefficient dropped} --- the picture predicts the table.

\textbf{The sawtooth is not noise.} $|x|$ is even and the odd $T_i$ are odd, so every odd
coefficient is exactly zero --- the orange line alternates between even coefficients and zeros.
\begin{center}\footnotesize
\begin{tabularx}{\linewidth}{@{}r >{\raggedleft\arraybackslash}X >{\raggedleft\arraybackslash}X@{}}
\toprule
$n$ & $e^{x}\sin 3x$ & $|x|$ \\
\midrule
$9$ & $10^{-3.6}$ & $10^{-1.2}$ \\
$17$ & $10^{-11.1}$ & $10^{-1.5}$ \\
$33$ & $\mathbf{10^{-14.3}}$ & $10^{-1.7}$ \\
$65$ & $10^{-13.8}$ & $10^{-2.0}$ \\
\bottomrule
\end{tabularx}
\par{\footnotesize max interpolation error}
\end{center}
\end{column}
\end{columns}
\vspace{1mm}
\src{Ledger \texttt{chebyshev}. The analytic case stops improving at $n=65$ because it has reached $10^{-14}$ --- rounding, not approximation.}
\end{frame}

%=============================================================================
\begin{frame}{More Than One Dimension: The Tensor Product}
\begin{columns}[T]
\begin{column}{0.55\textwidth}
\small
Chebyshev polynomials live on an interval; economic state spaces do not. The natural generalization multiplies them out. To approximate $F:[-1,1]^{d}\to\mathbb{R}$ at degree $\kappa$,
\[
  \hat{F}(x)=\sum_{n_1=0}^{\kappa}\cdots\sum_{n_d=0}^{\kappa}
  \theta_{n_1\dots n_d}\,T_{n_1}(x_1)\cdots T_{n_d}(x_d).
\]
\medskip
\textbf{The good news.} If the one-dimensional basis is orthogonal under a weight, the tensor-product basis is orthogonal under the product weight. Everything from the one-dimensional theory survives.
\medskip

\textbf{The bad news} is the number of sums.
\end{column}
\begin{column}{0.42\textwidth}
\begin{conceptbox}[title=\textbf{Count them}]\small
Each index runs over $\kappa+1$ values, and there are $d$ of them:
\[
  (\kappa+1)^{d}\ \text{coefficients.}
\]
\end{conceptbox}
\medskip
\begin{examplebox}\footnotesize
This is the same curse 3.A2 met with grids, arriving through a different door. Changing the representation did not remove it.
\end{examplebox}
\end{column}
\end{columns}
\vspace{1mm}
\src{Quant\_Macro Lec.~9, ``Tensors and Chebyshev Polynomial Approximation'' and ``Curse of Dimensionality in Tensor Approximation''.}
\end{frame}

%=============================================================================
\begin{frame}{The First Fix: Drop the High-Order Interactions}
\begin{columns}[T]
\begin{column}{0.52\textwidth}
\small
The tensor basis includes $x_1^{\kappa}x_2^{\kappa}\cdots x_d^{\kappa}$ --- a term of total degree $d\kappa$. Nobody needs that.
\medskip

\textbf{Complete polynomials} keep only the terms whose exponents sum to at most $\kappa$:
\[
  P^{d}_{\kappa}\equiv\Big\{x_1^{i_1}\cdots x_d^{i_d}\;\Big|\;\sum_{j=1}^{d}i_j\le \kappa\Big\}.
\]
Instead of letting every variable reach degree $\kappa$ independently, cap the \textbf{total} degree. Only low-complexity interactions survive --- which are the ones that matter for a smooth function.
\medskip

The count falls from $(\kappa+1)^{d}$ to $\binom{\kappa+d}{d}$.
\end{column}
\begin{column}{0.45\textwidth}
\begin{center}\footnotesize
\begin{tabularx}{\linewidth}{@{}r >{\raggedleft\arraybackslash}X >{\raggedleft\arraybackslash}X@{}}
\toprule
$d$ & tensor & complete \\
\midrule
$2$  & $25$ & $15$ \\
$3$  & $125$ & $35$ \\
$5$  & $3{,}125$ & $126$ \\
$10$ & $9{,}765{,}625$ & $\mathbf{1{,}001}$ \\
\bottomrule
\end{tabularx}
\end{center}
{\scriptsize\hfill $\kappa=4$}
\medskip
\begin{cautionbox}[title=\textbf{Still exponential?}]\footnotesize
No --- $\binom{\kappa+d}{d}$ is polynomial in $d$ for fixed $\kappa$. But it is still large, and it grows fast in $\kappa$.
\end{cautionbox}
\end{column}
\end{columns}
\vspace{1mm}
\src{Quant\_Macro Lec.~9, ``Complete Homogeneous Symmetric Polynomials''. \textbf{Gaspar \& Judd (1997)}, ``Solving Large-Scale Rational-Expectations Models,'' \emph{Macroeconomic Dynamics} 1(1), 45--75. Ledger \texttt{chebyshev.tensor\_vs\_complete}.}
\end{frame}

%=============================================================================
\begin{frame}{Which Terms Survive: $d=2$, $\kappa=2$}
\begin{columns}[T]
\begin{column}{0.56\textwidth}
\small
Two variables, degree two. Write each term as $x_1^{i_1}x_2^{i_2}$ and sort by \textbf{total degree}
$i_1+i_2$:
\medskip
\begin{center}\footnotesize
\begin{tabularx}{\linewidth}{@{}c >{\raggedright\arraybackslash}X c@{}}
\toprule
$i_1{+}i_2$ & \textbf{term} & \textbf{kept?} \\
\midrule
$0$ & $1$ & both \\
$1$ & $x_1,\;x_2$ & both \\
$2$ & $x_1^{2},\;x_1x_2,\;x_2^{2}$ & both \\
\midrule
$3$ & $x_1^{2}x_2,\;x_1x_2^{2}$ & tensor only \\
$4$ & $x_1^{2}x_2^{2}$ & tensor only \\
\bottomrule
\end{tabularx}
\end{center}
\medskip
The tensor rule caps each index separately and keeps all $3\times3=\mathbf{9}$. The complete rule
caps the \emph{sum} and keeps the first three rows --- $\mathbf{6}$.
\end{column}
\begin{column}{0.41\textwidth}
\begin{conceptbox}[title=\textbf{The rule in one line}]\small
Keep $x_1^{i_1}\cdots x_d^{i_d}$ if $\sum_j i_j\le\kappa$. Throw away everything else.
\end{conceptbox}
\medskip
\begin{examplebox}\footnotesize
Same rule at $d=3$, $\kappa=2$: the tensor basis keeps $3^{3}=27$ terms, the complete basis
$\binom{2+3}{3}=\mathbf{10}$. The gap widens with every dimension added.\\[1mm]
What you give up are the \textbf{high-order interactions} --- for a smooth function, the smallest
terms anyway, which is why the trade is usually free.
\end{examplebox}
\end{column}
\end{columns}
\vspace{1mm}
\src{Ledger \texttt{chebyshev.tensor\_vs\_complete}: $9\to6$ at $d=2$ and $27\to10$ at $d=3$, both with $\kappa=2$.}
\end{frame}

%=============================================================================
\begin{frame}{Smolyak Returns --- and It Is a Different Object Now}
\begin{columns}[T]
\begin{column}{0.55\textwidth}
\small
4.B1 built Smolyak sparse grids as a \textbf{node set}: where to evaluate a function you are interpolating, and where to put quadrature nodes. The construction was the nested ladder $m_1=1$, $m_i=2^{i-1}+1$ of Chebyshev extrema, with $G^1\subset G^2\subset\cdots$, combined over the index sets $\sum_j i_j\le d+\mu$.
\medskip

Here the \emph{same} index rule selects which \textbf{polynomials} enter the approximation:
\[
  \hat{F}(x)=\sum_{\mathbf{i}\,\in\,\Phi(q,d)}
  \theta_{\mathbf{i}}\;T_{i_1}(x_1)\cdots T_{i_d}(x_d),
\]
and the coefficients are chosen to kill a residual rather than to interpolate known values.
\end{column}
\begin{column}{0.42\textwidth}
\begin{conceptbox}[title=\textbf{Node set $\to$ basis}]\small
Same combinatorics, different job. There it answered \emph{where do I evaluate?} Here it answers \emph{which terms do I keep?}
\medskip
And the signed binomial weights of 4.B1 are what make the two views consistent.
\end{conceptbox}
\medskip
\begin{examplebox}\footnotesize
So a Smolyak collocation method uses the sparse set \emph{twice}: as the basis, and as the collocation points.
\end{examplebox}
\end{column}
\end{columns}
\vspace{1mm}
\src{Construction, nestedness and the index sets are \textbf{4.B1}'s --- see that deck rather than re-deriving them. Smolyak (1963); in economics Krueger \& Kubler (2004), Malin, Krueger \& Kubler (2011), Judd, Maliar, Maliar \& Valero (2014).}
\end{frame}

%=============================================================================
\begin{frame}{What It Buys, and What It Costs}
\begin{columns}[T]
\begin{column}{0.48\textwidth}
\small
4.B1's measured point counts, reused rather than recomputed:
\medskip
\begin{center}\footnotesize
\begin{tabularx}{\linewidth}{@{}r >{\raggedleft\arraybackslash}X >{\raggedleft\arraybackslash}X@{}}
\toprule
$d$ & Smolyak $\mu=2$ & tensor, $5$/dim \\
\midrule
$2$ & $13$ & $25$ \\
$5$ & $61$ & $3{,}125$ \\
$10$ & $\mathbf{221}$ & $9{,}765{,}625$ \\
\bottomrule
\end{tabularx}
\end{center}
\medskip
At $d=10$: \textbf{221 terms against nearly ten million}. And the run reproduces Malin, Krueger \& Kubler's published table exactly at $\mu=2$ --- $13,41,61,313$ for $d=2,4,5,12$.
\end{column}
\begin{column}{0.49\textwidth}
\begin{cautionbox}[title=\textbf{The curse is weakened, not broken}]\small
For $f\in C^{k}$ the error bound is
\[
  C_d\,n^{-k}(\log n)^{(d-1)(k+1)} ,
\]
the tensor rate spoiled by a logarithmic factor whose exponent \textbf{grows in $d$}. Sparsity postpones the curse; it does not repeal it.
\end{cautionbox}
\medskip
\begin{examplebox}\footnotesize
Which is the honest setup for S7: when $d$ is large enough that even this fails, the answer is a basis that \emph{adapts} rather than one chosen in advance.
\end{examplebox}
\end{column}
\end{columns}
\vspace{1mm}
\src{Ledger \texttt{smolyak} in \texttt{tools/figures/s04\_numbers.json} --- 4.B1's run, not recomputed here. Error bound: Barthelmann, Novak \& Ritter (1999).}
\end{frame}

%=============================================================================
\begin{frame}{Takeaways}
\begin{takeawaybox}
\begin{itemize}\small
  \item Chebyshev and Legendre are two members of the \textbf{Jacobi family}, distinguished by the orthogonality weight. Chebyshev's weight $1/\sqrt{1-x^{2}}$ emphasises the endpoints --- which is why its nodes cluster there.
  \item Bounded, smooth, stably generated by a three-term recursion, transformed by an FFT, and backed by error theorems. \textbf{When in doubt, use Chebyshev.}
  \item The interpolation theorem, in its sharp form: Chebyshev-node interpolation is within $\tfrac{2}{\pi}\log(n+1)+1$ of the \emph{best possible} polynomial of that degree --- a factor of $\mathbf{3.24}$ at $n=33$.
  \item Measured: for an analytic function the error reaches \textbf{machine precision by degree 33}; for $|x|$ it falls like $1/n$ and sixty-five terms buy two digits. \textbf{The basis is not the problem --- the function is.}
  \item The tensor product preserves orthogonality and costs $(\kappa+1)^{d}$ coefficients. \textbf{Complete polynomials} cut that to $\binom{\kappa+d}{d}$: $1{,}001$ instead of $9.8$ million at $d=10$, $\kappa=4$.
  \item \textbf{Smolyak}, 4.B1's node set, is here a \textbf{basis}: the same index rule chooses which polynomials enter. $221$ terms at $d=10$ against nearly ten million --- and an error bound that weakens the curse without breaking it.
\end{itemize}
\end{takeawaybox}
\end{frame}

%=============================================================================
\begin{frame}{Bridge: When the Function Has a Corner}
\begin{columns}[T]
\begin{column}{0.53\textwidth}
\small
Everything in this deck rewards \textbf{smoothness}, and the measured penalty for its absence was severe: $|x|$ needed sixty-five Chebyshev terms to reach two digits, and 5.B2 showed the oscillation beside a kink never goes away.
\medskip

Economics supplies corners as a matter of course --- borrowing limits, irreversibility, non-negativity. So there has to be a basis that does not care.
\end{column}
\begin{column}{0.44\textwidth}
\begin{conceptbox}[title=\textbf{Next (5.B4)}]\small
Finite elements: a \textbf{local} basis. Chop the domain into pieces, put a simple function on each, and let a piece boundary sit exactly where the kink is. Slow to converge, indifferent to smoothness, and sparse.
\end{conceptbox}
\end{column}
\end{columns}
\end{frame}

\end{document}
